\section{Introduction}
\label{sec:intro}

Language-model agents increasingly act on external systems: they publish posts, charge cards, send email,
open tickets and trigger deployments through tool calls. Those calls cross a network, and networks lose
messages. When a write times out or returns a server error, the action may or may not have taken effect.
A client that retries blindly can charge a customer twice or announce a release twice; a client that gives
up can silently skip required work. Distributed systems treat this ambiguity as fundamental---a client
cannot distinguish a lost request from a lost acknowledgement---and have converged on remedies that live in
the \emph{interface} rather than in the client: idempotency keys, conditional writes and queryable operation
status \citep{helland2012idempotence,stripe_idempotent,featonby2020retries}.

Agents inherit the ambiguity, but not by default the remedies. Concurrent work has begun to document the
consequences. \citet{sun2026didithappen} constructs observation-equivalent counterfactual pairs---identical
timeouts with and without a hidden commit---and finds that five open and API models retry blindly about half
the time unless they are given a status query or a stable key. IdempotencyBench
\citep{gopnalswamy2026idempotencybench} argues from scripted agents that exactly-once behaviour is a property
of the execution substrate rather than of prompting, and \citet{mansoor2026verified} show that a
verify-then-retry wrapper around a single model lowers duplicate rates. These studies leave the central
engineering question open: \emph{where should exactly-once be enforced?} If frontier models reliably verify
before they retry, better models suffice. If agent harnesses---the command-line agents and frameworks that
wrap a model---retry transparently or change what the model sees, the harness matters. If some failures cannot
be resolved by any amount of verification, the remedy must be in the tool contract.

We study this question with \bench{}, a benchmark in which every recovery decision is graded against a
ground-truth ledger of committed side effects, and with a factorial experiment that varies the model, the
harness, the tool contract and the recovery policy while holding the task, the world and the fault fixed.
\bench{} injects twelve fault modes at the service boundary, including three that prior agent benchmarks
omit: \emph{late commits}, where a timed-out request is still in flight and lands after the client has
moved on (with a fixed or a heavy-tailed delay); \emph{redelivery}, where an at-least-once transport delivers a
request twice; and \emph{partial batches}. Its tools carry realistic contracts---optional idempotency keys,
eventually consistent read paths with documented lag, endpoints with no read-back at all, naturally idempotent
operations---so that the same fault can be recoverable, recoverable only with care, or unrecoverable depending on
the contract. Because the sandbox is exposed through the Model Context Protocol, the identical environment runs
under a minimal function-calling scaffold and under production agent harnesses.

Our study covers \NModels{} recent models from four providers in a minimal scaffold and three production
harnesses (GitHub Copilot CLI, Hermes and Codex CLI), for a total of \NEpisodesAll{} graded episodes.
The answer to our question depends on the kind of failure, and we find:

\begin{itemize}[leftmargin=1.2em,itemsep=2pt]
\item \textbf{When a read-back can resolve the failure, the model decides.} Frontier models instructed to
act exactly once almost never duplicate a write whose acknowledgement was lost (\FrontierPostDSR{}\% of
episodes), because they verify before retrying; a server error that hides a commit is harder (\FrontierHttpPostDSR{}\%
for frontier models, \WeakHttpPostDSR{}\% for the weaker \WeakList{}), and without the
explicit instruction the weaker and faster models duplicate markedly more while the strongest are unaffected.
On these faults the model explains \ShareModelOneRes{}\% of the explained variance in duplicates.
\item \textbf{When it cannot, only the contract does.} When the timed-out request is still in flight, the
same verification finds nothing and the retry creates a duplicate in \FrontierLateDSR{}\% of frontier
episodes; under redelivery, in \FrontierRedelivDSR{}\%. Here the contract explains \ShareContractOneUnres{}\% of
the variance. We prove that no verification-only policy can be exactly-once under late commits without a known
bound on in-flight time, and measure the price of waiting for such a bound: waiting works when the bound is short
and known, but with heavy-tailed in-flight delays even an hour of waiting per episode falls short of what keys
achieve with no added latency. Offering keys is enough: under a counterfactual contract in which every write accepts
one, agents attach them in \EKeysVanillaKeyUse{}\% of episodes and the duplicate rate falls from \ETwoVanillaDSR{}\%
to \EKeysVanillaDSR{}\% with no other change; every remaining duplicate involves an agent that sent the first attempt
without a key.
\item \textbf{Harnesses matter little; client-side middleware has a ceiling.} Three production harnesses and a
minimal scaffold running the same model behave almost identically. \GuardSummary{}
\item \textbf{Agents do not know when they have duplicated.} In \OverclaimGivenDup{}\% of episodes that
produced a duplicate, the agent reported the task as completed.
\end{itemize}

We release the benchmark, the harness adapters and every episode trace\CodeCite.\footnote{Code (MIT licence) and
data (CC BY 4.0): \CodeURL.}
